\cotefile{15}
\scanpage{207}
\uline{Faisons varier} $x$ pt complexe (ou \unsure{pl\textsuperscript{ce}|place|plus} archimédienne). Je pense que \uline{les $\unsure{\xi|\zeta}_{x,\mathbb Z}$ et $=$ les $\xi_{x,\mathbb R}$} \inserted{$[\simeq \xi_{x,\mathbb C}$ muni de $\tau]$} \uline{qui correspondent à des \unsure{pl\textsuperscript{ces}|places} distinctes sont non isomorphes}
% slanted
\margin{\tikz[baseline=(a.base)]\node[draw,ellipse,inner sep=1pt,align=center](a){faux, cf\\groupoïde\\fondamental\\sur $Z$};}

Quant aux $\xi_{x,\mathbb C}=\xi_{x,\infty}$, eux $=$s, \deleted{\unsure{comme|isomorphes}} \inserted{\unsure{comme}} $\unsure{K|k}$ est \deleted{nous avons dit, il sont} alg. clos ils sont isomorphes entre eux, mais pas canoniquement ....

\uline{NB} En car $p>0$, je pense que $T_{\alpha}(\mathbb R)=\varnothing$ pour $\alpha$ grand, ce qui explique qu'il n'y ait pas de phénomènes de ``composantes connexes'' à considérer alors .... ``Ramification en $\mathbb R$'' \ldots
% slanted
\margin{vérifier si il ne peut y avoir des pt à val. dans une ext. tot. réelle ???}

Si

\begin{center}\rule{3cm}{0.5pt}\\[-0.65em]\rule{2.6cm}{0.5pt}\end{center}

Pt motivique de $X$ à valeurs dans un corps (ou un anneau) $k$ : Objets de $T(k)$. Langage de la \uline{fibre d'un motif en un point}

les p\textsuperscript{ts} forment une catégorie fibrée sur $(Sch)$ muni de fpqc, \deleted{\illegible{1}} c'est un \uline{torseur}. Pour $k$ donné, les isom. d'un \deleted{objet} $\xi\in T(k)$ sur un autre \deleted{\unsure{forment|jouent}} les
% archive sheet number: 206

\scanpage{208}
SYSTÈMES de RACINES \mspage{(3)}\hfill Tribu n\textsuperscript{o} 59

linéaire réelle $\alpha\in\underline{\mathrm a}^{*}$ sur $\underline{\mathrm a}$, soit $\underline{\mathrm g}^{\alpha}=\{\mathrm x\in\underline{\mathrm g}\mid[\mathrm a,\mathrm x]=(\mathrm a).\mathrm x\;(\mathrm a\in\underline{\mathrm a})\}$. Alors l'ensemble $\mathrm R$ des éléments $\alpha\in\underline{\mathrm a}^{*}-0$ pour lesquels $\underline{\mathrm g}^{\alpha}\ne(0)$ est un système de racines, en général non réduit. Le groupe de Weyl de $\mathrm R$ s'identifie au groupe des automorphismes de $\underline{\mathrm a}$ induits par des éléments de $\mathrm{Ad}_{\underline{\mathrm g}}\,\underline{\mathrm k}$ laissant $\underline{\mathrm a}$ stable.

III. Soient $\mathrm k$ un corps, $\mathrm G$ un groupe linéaire algébrique connexe semi-simple sur $\mathrm k$, $\mathrm T$ un tore algébrique sur $\mathrm G$ décomposé sur $\mathrm k$ maximal, et $\mathrm X^{*}$ le groupe des caractères rationnels de $\mathrm T$. Le dual $\mathrm X_{*}$ de $\mathrm X^{*}$ s'identifie canoniquement au groupe des homomorphismes rationnels du groupe multiplicatif $\mathrm G_{\mathrm m}$ dans $\mathrm T$. Soit $\underline{\mathrm g}$ l'algèbre de Lie de $\mathrm G$ et, pour $\alpha\in\mathrm X^{*}$, soit $\underline{\mathrm g}^{\alpha}=\{\mathrm x\in\underline{\mathrm g}\mid\mathrm{Ad}\,\mathrm t(\mathrm x)=\alpha(\mathrm t).\mathrm x\;(\mathrm t\in\mathrm T)\}$. Alors l'ensemble $\mathrm R$ des éléments non nuls $\alpha\in\mathrm X^{*}$ pour lesquels $\mathrm g^{\alpha}\ne(0)$ est un système de racines, qui est réduit si $\mathrm T$ est aussi un sous-groupe de Cartan de $\mathrm G$, ce qui a en particulier toujours lieu lorsque $\mathrm k$ est algébriquement clos. Le groupe de Weyl $\mathrm W(\mathrm R)$ s'identifie au quotient du normalisateur de $\mathrm T$ par le centralisateur de $\mathrm T$ dans $\mathrm G$.

Si $\mathrm k=\underline{\underline{\mathrm C}}$ l'algèbre de Lie $\underline{\mathrm t}$ de $\mathrm T$ est une sous-algèbre de Cartan de $\underline{\mathrm g}$ et $\mathrm R$ est isomorphe au système de racines défini en I. Si $\mathrm k=\underline{\underline{\mathrm R}}$, alors l'ensemble des points réels de $\underline{\mathrm t}$ est conjugué à l'algèbre $\underline{\mathrm a}$ considérée en II, et $\mathrm R$ est isomorphe au système de racines défini en II.

IV. Soient $\mathrm G$ un groupe de Lie compact connexe semi-simple, $\underline{\mathrm g}$ son algèbre de Lie, $\mathrm T$ un tore maximal de $\mathrm G$ et $\mathrm X$ le groupe des homomorphismes continus de $\mathrm T$ dans $\underline{\underline{\mathrm C}}^{*}$. Pour tout $\alpha\in\mathrm X$, soit $\underline{\mathrm g}^{\alpha}=\{\mathrm x\in\underline{\mathrm g}\mid\mathrm{Ad}\,\mathrm t(\mathrm x)=\alpha(\mathrm t).\mathrm x\;(\mathrm t\in\mathrm T)\}$. Alors l'ensemble $\mathrm R$ des éléments $\alpha\in\mathrm X-0$ pour lesquels $\underline{\mathrm g}^{\alpha}\ne(0)$ est un système de racines réduit dans $\mathrm V=\mathrm X\otimes\underline{\underline{\mathrm Q}}$, dont le groupe de Weyl s'identifie au quotient par $\mathrm T$ du normalisateur de $\mathrm T$ dans $\mathrm G$. On a $\mathrm Q(\mathrm R)\subset\mathrm X\subset\mathrm P(\mathrm R)$ et $\mathrm X=\mathrm P(\mathrm R)$, (resp. $\mathrm X=\mathrm Q(\mathrm R)$), si $\mathrm G$ est simplement connexe, (resp. isomorphe à $\mathrm{Ad}\,\underline{\mathrm g}$). L'algèbre de
% archive sheet number: 207

\scanpage{209}
rôle de \uline{classes de chemins}, ils sont déterminés \inserted{(s'il y en a au moins un)} \unsure{mod} automorphismes \unsure{du premier|des premiers}, qui font les p\textsuperscript{ts} d'un groupe algébrique. Mais en général il peut y avoir des p\textsuperscript{ts} à valeurs dans un $=$ $k$ non isomorphes. Les groupes de Galois \unsure{opèrent|opérant|opère} sur \deleted{les points} \unsure{l'en\textsuperscript{s}|l'un|l'ens.} des p\textsuperscript{ts} à valeurs dans $k$. Si on a un pt $\unsure{\xi|\zeta}$ à val \unsure{dans} $k$, tous les autres p\textsuperscript{ts} s'en déduisent par des fibrés principaux \unsure{hom.} sur le (pro)groupe algébrique $\underline{\mathrm{Aut}}(\xi)$.

\diagram{\[
\begin{array}{@{}l@{}}
\left\{\begin{array}{ll}
\text{Pts $\ell$-adiques} & \text{valeurs }\mathbb Z_{\ell}\text{ ou }\mathbb Q_{\ell}\\
\text{pts transcendants} & \text{valeurs }\mathbb R\ \deleted{\text{ou }\mathbb C}\ \deleted{\unsure{\text{déjà}}}\text{, voire }\mathbb Z\\
 & \text{[structure supplémentaire]}\\
 & \text{Il semble que ces p\textsuperscript{ts} correspondent}\\
 & \text{\unsure{biunivoq\textsuperscript{mt}|biunivoquent|biunivoquement} aux pts de}\\
 & X\text{ à valeurs dans }\mathbb C\text{,}\\
 & \text{mod conjugaison}
\end{array}\right.\\[0.5em]
\begin{array}{ll}
\text{pts Hodge-de Rham} & \text{\unsure{valeur|valeurs} }k\text{ \unsure{quelconques|quelconque}}\\
\text{pts de Witt} & \text{valeurs }W_{p}(k)\ \inserted{\text{$k$ corps de car $p$}}
\end{array}
\end{array}
\]}
% archive sheet number: 208

\scanpage{210}
SYSTÈMES de RACINES \mspage{(2)}\hfill Tribu n\textsuperscript{o} 59

Pour orienter le lecteur, nous indiquons ici les principaux exemples de systèmes de racines que l'on rencontre dans la théorie des groupes \unsure{$\cdot$}et algèbres de Lie.

I. Soient $\underline{\mathrm g}$ une algèbre de Lie semi-simple sur $\underline{\underline{\mathrm C}}$ et $\underline{\mathrm h}$ une sous-algèbre de Cartan de $\underline{\mathrm g}$. L'algèbre $\underline{\mathrm h}$ est donc commutative, formée d'éléments semi-simples. Une racine de $\underline{\mathrm g}$ par rapport à $\underline{\mathrm h}$ est une forme linéaire $\alpha$ non nulle sur $\underline{\mathrm h}$ pour laquelle il existe $\mathrm x\in\underline{\mathrm g}$, $\mathrm x\ne0$, tel que $[\mathrm h,\mathrm x]=\alpha(\mathrm x)$ quel que soit $\mathrm h\in\underline{\mathrm h}$. Soit $\mathrm V$ l'espace vectoriel sur $\underline{\underline{\mathrm Q}}$ engendré par l'ensemble $\mathrm R$ des racines de $\underline{\mathrm g}$ par rapport à $\underline{\mathrm h}$. Alors $\underline{\mathrm h}^{*}=\mathrm V\otimes\underline{\underline{\mathrm C}}$, et $\mathrm R$ est un système de racines réduit dans $\mathrm V$, au sens du \S~6. Tout système de racines réduit est l'ensemble des racines d'une algèbre semi-simple complexe, unique à son isomorphisme près. Tout automorphisme de $\underline{\mathrm g}$ laissant stable $\underline{\mathrm h}$ induit un automorphisme de $\underline{\mathrm h}^{*}$ laissant $\mathrm R$ et $\mathrm V$ invariants, et le groupe des automorphismes de $\mathrm R$ ainsi obtenus à partir des automorphismes, (resp. des automorphismes intérieurs), de $\underline{\mathrm g}$ s'identifie au groupe de tous les automorphismes de $\mathrm R$, (resp. au groupe de Weyl de $\mathrm R$)\unsure{:|.} Identifions le dual de $\mathrm V$ à une $\underline{\underline{\mathrm Q}}$-forme $\underline{\mathrm h}^{\unsure{\mathrm o|0|\circ}}$ de $\underline{\mathrm h}$ et soit $\mathrm G$ un groupe de Lie complexe connexe d'algèbres de Lie $\underline{\mathrm g}$. Alors le noyau $\Gamma$ de l'application exponentielle de $\underline{\mathrm h}$ dans $\mathrm G$ est un réseau de $\mathrm i.\underline{\mathrm h}^{\mathrm o}$ compris entre $2\pi\mathrm i.\mathrm Q(\mathrm R^{*})$ et $2\pi\mathrm i.\mathrm P(\mathrm R^{*})$, où $\mathrm R^{*}$ est le système de racines inverse de $\mathrm R$, qui est égal à $2\pi\mathrm i.\mathrm Q(\mathrm R^{*})$, (resp. $2\pi\mathrm i\mathrm P(\mathrm R^{*})$), si $\mathrm G$ est simplement connexe, (resp. isomorphe à $\mathrm{Ad}\,\underline{\mathrm g}$). Le groupe $2\pi\mathrm i.\mathrm P(\mathrm R^{*})/\Gamma$ s'identifie au centre de $\mathrm G$, et $\Gamma/2\pi\mathrm i.\mathrm Q(\mathrm R^{*})$ au groupe fondamental de $\mathrm G$. Le réseau de $\mathrm V$ associé à $(2\pi\mathrm i)^{-1}.\Gamma$ est engendré par les poids des représentations de $\mathrm G$.

II. Soient $\underline{\mathrm g}$ une algèbre de Lie semi-simple sur $\underline{\underline{\mathrm R}}$, $\underline{\mathrm k}$ une sous-algèbre compacte maximale de $\underline{\mathrm g}$, $\underline{\mathrm p}$ le complément orthogonal de $\underline{\mathrm k}$ par rapport à la forme de Killing, et $\underline{\mathrm a}$ une sous-algèbre commutative maximale de $\underline{\mathrm p}$. Pour toute forme
% archive sheet number: 209

\scanpage{211}
\uline{\unsure{Question|Questions}}
% slanted
\margin{\tikz[baseline=(a.base)]{\draw[gray] (.1,-1.9) .. controls (.25,-.4) and (.17,1.2) .. (.15,2.4); \node[anchor=east,rotate=35](a) at (0,0){À \unsure{revoir|prévoir|voir}};}}

\diagram{\[
\left\{\begin{array}{@{}l@{}}
\left\{\begin{array}{@{}l@{}}
\parbox{.79\linewidth}{pts \inserted{de $X$} à val. dans $k$ \inserted{de car $p>0$} : correspondent aux pt motiviques \inserted{dans $W(k)$ (ou corps des fractions ?)} avec \uline{structure supplémentaire} \tikz[baseline=(a.base)]\node[ellipse,draw,inner sep=1pt](a){$F\deleted{\unsure{F|S}},V$}; [foncteurs des groupes formels]}\\
\begin{tikzcd}[ampersand replacement=\&,column sep=.8em,row sep=.4em]
{} \arrow[dr,squiggly] \& {}\\
{} \& \parbox{.60\linewidth}{pts motiviques généralisés (au \unsure{sens|une} cat. dérivées) dans $k$, avec structure filtrée !}
\end{tikzcd}
\end{array}\right.\\[.6em]
\parbox{.82\linewidth}{pts \inserted{de $X$} à val. dans $k$ de car $0$ : correspondent aux pts motiviques dans $k$, avec \uline{structure filtrée}}
\end{array}\right.
\]}

\diagram{\[
\left|\begin{array}{l}
\text{pts \inserted{de $X$} à valeurs dans $\mathbb C$ mod conj. complexe}\\
\text{(= \unsure{places|pl\textsuperscript{ces}|plus})}\ \longleftrightarrow\
 \text{pts \inserted{motiviques} à valeurs dans }
 \tikz[baseline=(a.base)]\node[draw,ellipse](a){$\mathbb R$};\text{, ou}\\
\quad\text{à valeurs dans }\tikz[baseline=(a.base)]\node[draw,ellipse](a){$\mathbb C$ avec \tikz[baseline=(b.base)]\node[draw,ellipse,inner sep=1pt](b){$f_{\infty}$};};\text{, ou}\\
\quad\text{à valeurs dans }\mathbb Z
\end{array}\right.
\]}

\uline{\deleted{\unsure{Sch de car $p$|sur $X$ de car $p$|de $X$ de car $p$|Sch de cas $p$}}}

\uline{Pts banaux}\qquad Pts à valeurs dans \tikz[baseline=(a.base)]\node[draw,ellipse](a){$\mathbb Z_{\unsure{l|\ell}}$};, $\ell$ premier aux \unsure{car.} résiduelles. ou \tikz[baseline=(a.base)]\node[draw,ellipse](a){$\mathbb Q_{\ell}$};

! \uline{Cas particulier important}\qquad \uline{Pts à valeurs dans corps} fini $\mathbb F_q$ correspondent aux p\textsuperscript{ts} à valeurs dans les \tikz[baseline=(a.base)]\node[draw,ellipse](a){$W(\mathbb F_q)$};, avec structure supplémentaire $(F,V)$
% archive sheet number: 210

\scanpage{212}
\hfill Tribu n\textsuperscript{o} 59

\begin{center}
SYSTÈMES de RACINES

\uline{Introduction}, (Rédaction Borel).
\end{center}

Ce chapitre est principalement consacré à l'étude de certains ensembles finis de vecteurs d'un espace euclidien, appelés systèmes de racines, et des groupes de déplacements d'un espace euclidien vectoriel ou affine engendré par des réflexions. Son but est de fournir le cadre naturel dans lequel s'expriment et se démontrent un grand nombre de propriétés importantes des groupes et algèbres de Lie semi-simples sur $\underline{\underline{\mathrm R}}$ ou sur $\underline{\underline{\mathrm C}}$ et des groupes algébriques semi-simples sur un corps quelconque. Bien que la rédaction ait été dictée par les besoins des chapitres suivants de ce Livre, nous avons été amenés en plusieurs endroits à élargir de diverses façons les hypothèses faites sur les groupes envisagés, de manière à placer certains résultats dans leur domaine naturel de validité, qui déborde parfois celui des applications que nous avons en vue. Le lecteur uniquement intéressé à ces dernières peut concentrer son attention sur les systèmes de racines, (\S~6), et sur leurs groupes de Weyl vectoriels ou affines. Il peut se limiter, dans le \S~1 aux réflexions d'un espace euclidien ou d'un $\underline{\underline{\mathrm Z}}$-module libre, dans le Théorème 1 du \S~7 à l'implication $\mathrm{(a)}\Longrightarrow\mathrm{(b)}$ lorsque $\mathrm k=\underline{\underline{\mathrm R}}$ et que $\mathrm G$ est engendré par des réflexions, et envisager la discussion des groupes de Coxeter faite au \S~4 avant tout comme un auxiliaire permettant de donner une démonstration unifiée des propriétés fondamentales des groupes de déplacements linéaires ou affines engendrés par des réflexions\textsuperscript{*}.

\noindent\makebox[\linewidth][s]{- - - - - - - - - - - - - - - - - - - - - - - - - - - - - - - - - - - - - - - - -}

($^{*}$) Pour rédiger ce paragraphe, nous avons largement utilisé des Notes de Séminaire, (non publiées), de Jacques Tits.
% archive sheet number: 211

\scanpage{213}
\uline{Groupe fondamental \unsure{motifique|motivique} discret}

Soit $\unsure{\xi|\zeta}\in T(\mathbb Z)$, on pose
\[
\begin{aligned}
\pi_1(X,\xi)&=\mathrm{Aut}(\xi).\\
             &=G_{\xi}(\mathbb Z),
\end{aligned}
\]

\uline{N.B.} En vertu de Borel\inserted{\unsure{-Harish|-Harsh}}, il n'y a qu'un nb fini de $\xi\in T_{\alpha}(\mathbb Z)$ possibles, pour $\alpha$ donné. D'où un nb fini de \uline{groupes fondamentaux discrets}. Sont ils tous associés à des \deleted{\unsure{deux|choses}} places ??

\[
\begin{aligned}
&\xi\in T_{(\alpha)}(k),\ k\text{ un anneau}\\
&M\text{ motif }\in\mathrm{Ob}\,\unsure{\mathscr M|\mathcal M}_{\alpha}\\
&M_{\xi}\text{ fibre en }\xi,\qquad M_{\xi}=\xi(M)\\
&\pi_1(X,\xi)_{\deleted{\unsure{\alpha|\mathrm{an}|0}}}=\mathrm{Aut}(\xi),\quad
G(X,\xi)=\underline{\mathrm{Aut}}(\xi)\\
&\deleted{\unsure{G_{\xi_{\!k}}|G_{\xi}}}\qquad \pi_1(X,\xi)_{\unsure{\alpha|0|\mathrm{an}}}\simeq G(X,\xi)(k)
\end{aligned}
\]

Comparaison avec $\pi_1(X^{\unsure{\mathrm{top}|\mathrm{an}}},\unsure{\xi|a})$ quand $\xi$ associé à une \uline{place} on a ....

\uline{Ex}. $X$ connexe sur $\mathbb C$, $\xi$ défini par un pt de $X$ à valeurs dans $\mathbb C$. Le foncteur $\xi_{\mathbb C}^{\deleted{\unsure{\mathbb C|a}}}$ a une structure \uline{filtrée}. $\xi$ considéré comme foncteur à valeurs dans les modules sous $\pi_1(X,\xi)$ munis d'une \uline{structure de Hodge} est \unsure{pl|pl.|pl\textsuperscript{t}} fidèle ....
% archive sheet number: 212

\scanpage{214}
\mspage{1 - 49 -}\hfill n\textsuperscript{o} 413

Soient $\mathrm G$ et $\mathrm H$ deux noyaux de groupe de classe $\mathrm C^{\mathrm r}$. Un \uline{morphisme} $\psi$ de $\mathrm G$ dans $\mathrm H$ est une application de classe $\mathrm C^{\mathrm r}$ de $\mathrm G$ dans $\mathrm H$ satisfaisant aux conditions suivantes :

(a) $\psi(\mathrm e_{\mathrm G})=\mathrm e_{\mathrm H}\qquad\psi(\unsure{\vartheta|\theta}_{\mathrm G}(\mathrm x))=\vartheta_{\mathrm H}(\psi(\mathrm x))$ pour tout $\mathrm x\in\mathrm G$ ;

(b) si $(\mathrm x,\mathrm y)\in\Omega(\mathrm G)$, alors $(\psi(\mathrm x),\psi(\mathrm y))\in\Omega(\mathrm H)$ et on a :
\[
\mathrm m_{\mathrm H}(\psi(\mathrm x),\psi(\mathrm y))=\psi(\mathrm m_{\mathrm G}(\mathrm x,\mathrm y))
\]
(autrement dit $\psi$ est une représentation pour les lois de composition non partout définies sur $\mathrm G$ et $\mathrm H$). Nous laisserons au lecteur le soin d'étendre au cas des noyaux de groupe les propositions 3 et 3 bis.

Soit $\mathrm G$ un noyau de groupe. Il existe un voisinage ouvert $\mathrm V$ de $\mathrm e$ tel que les conditions de (iii) soient réalisées pour $\mathrm x,\mathrm y,\mathrm z\in\mathrm V$. En particulier le \uline{produit de convolution} associé à la restriction de $\mathrm m$ à $\mathrm V\times\mathrm V$ définit une loi de composition sur $\mathrm T_{\mathrm e}^{(\mathrm r)}(\mathrm G)$, qui est bilinéaire et associative (bien que non partout définie si $\mathrm r<\infty$). Nous désignerons encore par $\mathrm U(\mathrm G)$ l'espace $\mathrm T_{\mathrm e}^{(\mathrm r)}(\mathrm G)$ muni de cette multiplication. Nous laisserons au lecteur le soin d'étendre les autres notations et résultats des n\textsuperscript{os} 12 à 16. Signalons par exemple que si $\mathrm f$ est un morphisme de $\mathrm G$ dans un noyau de groupe $\mathrm H$, alors $\mathrm f_{*}$ est une \uline{représentation} de $\mathrm U(\mathrm G)$ dans $\mathrm U(\mathrm H)$.

De plus, l'application tangente $\mathrm m^{\mathrm T}$ à $\mathrm m$ est une application de classe $\mathrm C^{\mathrm r-1}$ de $\mathrm T(\quad)$ dans $\mathrm T(\mathrm G)$. On peut identifier $\mathrm T(\quad)$ à un ouvert de $\mathrm T(\mathrm G)\quad\mathrm T(\mathrm G)$ et on vérifie aisément que $\mathrm T(\mathrm G)$ muni du point marqué $0_{\mathrm e}$, de l'application $\mathrm m^{\mathrm T}$ de $\mathrm T(\quad)$ dans $\mathrm T(\mathrm G)$ et de l'application $\vartheta^{\mathrm T}$ de $\mathrm T(\mathrm G)$ dans $\mathrm T(\mathrm G)$ est un \uline{noyau de groupe} de classe $\mathrm C^{\mathrm r-1}$ ; la projection canonique est un morphisme de classe $\mathrm C^{\mathrm r-1}$ de $\mathrm T(\mathrm G)$ sur $\mathrm G$. Nous laisserons au lecteur le soin d'étendre au cas des noyaux de groupe les propositions 11 et 12 ainsi que leurs corollaires.
% archive sheet number: 213

\scanpage{215}
(Résulte de la conjonction des conj de Hodge et de Tate \inserted{\tikz[baseline=(a.base)]\node[draw,ellipse,inner sep=1pt](a){sous forme faible.};} \unsure{?}, du moins si $X$ \uline{\unsure{normal|connexe|réduit}}). \uline{\uuline{\deleted{\unsure{Caractérisation}}Caractériser \deleted{de} l'image}} ??
\margin{$\vert\,\vert\,\vert\,\vert\,\vert$}

M conclusion si $X$ défini sur $k\subset\mathbb C$, $k$ alg. clos, p. ex. $k=\overline{\mathbb Q}$.

\begin{center}\uline{Analyser l'exemple de Serre de} $\pi_1$ \uline{non isomorphes}\end{center}

[\uline{Cas corps de base} $\mathbb R$ : tout pt \deleted{réel} \inserted{fermé $a$} définit une $\unsure{\xi|\zeta}_a\in T_{\unsure{\alpha|a}}(\mathbb Z)$. Tous ces $\xi_a$ sont isomorphes si \deleted{\unsure{a}} $X$ est \deleted{\unsure{géom.} \uline{connexe}, \unsure{sinon}} \inserted{\unsure{irrid|irréd} (pas néc. géom \unsure{irrid|irréd})} \deleted{il peut y en avoir deux non isom.}

Si $a$ est un pt réel, alors $\xi_a$ admet un \unsure{autom} $\mathrm{frob}_{\infty}$ d'ordre $2$, qui avec la comp. conn. de $G(X,\xi)$ engendre $G(X,\xi)$.
\margin{\diagram{\begin{tikzcd}[ampersand replacement=\&,row sep=.5em,column sep=.6em]
{} \& X\arrow[dd,dash]\\
\mathbb C\arrow[ur,dash]\arrow[dr,dash] \& {}\\
{} \& \mathbb R
\end{tikzcd}}}
Le foncteur
\[
\begin{aligned}
M\ \unsure{\sim|\leadsto}\ (\unsure{\xi(M_{\xi})|\xi(M)_{\xi}|\xi(M_{\mathbb C})},\ \unsure{+^{\prime}|+}&\text{structure de Hodge}+\text{\unsure{réellification|réalification} définie}\\
&\text{par }\mathrm{frob}_{\infty}\,)
\end{aligned}
\]
est pleinement fidèle. ]

\margin{$\vert\,\vert\,\vert$}
\uline{Pb} étudier comment varie la structure de Hodge pour un pt \inserted{complexe} variable, \uline{en fonction de la donnée \unsure{départ|discret|discrète|définie}} en le pt $a_0$,
% archive sheet number: 214

\scanpage{216}
\mspage{1 - 48 -}\hfill n\textsuperscript{o} 413

17. \uline{Noyaux de groupe}.

DÉFINITION 5.- \uline{On appelle noyau de classe} $\mathrm C^{\mathrm r}$ \uline{une variété} $\mathrm V$ \uline{de classe} $\mathrm C^{\mathrm r}$ \uline{munie d'un point marqué} $\mathrm e$. \uline{Un noyau de groupe de classe} $\mathrm C^{\mathrm r}$ \uline{est un noyau} $\mathrm G$ \uline{muni d'une application} $\unsure{\vartheta|\theta}_{\mathrm G}$ \uline{de classe} $\mathrm C^{\mathrm r}$ \uline{de} $\mathrm G$ \uline{dans} $\mathrm G$ \uline{et d'une application} $\mathrm m_{\mathrm G}$ \uline{de classe} $\mathrm C^{\mathrm r}$ \uline{d'un ouvert} $\Omega(\mathrm G)$ \uline{de} $\mathrm G\times\mathrm G$ \uline{dans} $\mathrm G$ \uline{satisfaisant aux conditions suivantes} :

\begingroup
\makeatletter\tagsleft@true\makeatother
(i) $\Omega$ \uline{contient les points} $(\mathrm x,\vartheta(\mathrm x))$ \uline{et} $(\vartheta(\mathrm x),\mathrm x)$ \uline{et on a}
\[
\mathrm m(\mathrm x,\vartheta(\mathrm x))=\mathrm m(\vartheta(\mathrm x),\mathrm x)=\mathrm e \tag*{(36)}
\]
\uline{pour tout} $\mathrm x\in\mathrm G$ ;

(ii) $\Omega$ \uline{contient les points} $(\mathrm e,\mathrm x)$ \uline{et} $(\mathrm x,\mathrm e)$ \uline{et on a}
\margin{\handwritten{\unsure{$\parallel$|T|$\pi$}}}
\[
\mathrm m(\mathrm e,\mathrm x)=\mathrm m(\mathrm x,\mathrm e)=\mathrm x \tag*{(37)}
\]
\uline{pour tout} $\mathrm x\in\mathrm G$ ;

(iii) \uline{si} $\mathrm x,\mathrm y,\mathrm z$ \uline{sont trois éléments de} $\mathrm G$ \uline{tels que} $(\mathrm x,\mathrm y)$, $(\mathrm y,\mathrm z)$, $(\mathrm m(\mathrm x,\mathrm y),\mathrm z)$ \uline{et} $\mathrm m(\mathrm x,\mathrm m(\mathrm y,\mathrm z))$ \uline{appartiennent à} $\Omega$\handwritten{$\unsure{\searrow|,}$} \uline{on a} :
\[
\mathrm m(\mathrm m(\mathrm x,\mathrm y),\mathrm z)=\mathrm m(\mathrm x,\mathrm m(\mathrm y,\mathrm z)). \tag*{(38)}
\]
\endgroup

\[
\begin{aligned}
\text{On a }\mathrm e=\mathrm m(\mathrm e,\vartheta(\mathrm e))=\vartheta(\mathrm e)
\quad\text{et }\mathrm x&=\mathrm m(\mathrm e,\mathrm x)=\mathrm m(\mathrm m(\vartheta(\vartheta(\mathrm x)),\vartheta(\mathrm x)),\mathrm x)=\\
&=\mathrm m(\vartheta(\vartheta(\mathrm x)),\mathrm m(\vartheta(\mathrm x),\mathrm x)=\mathrm m(\vartheta(\vartheta(\mathrm x)),\mathrm e)=\vartheta(\vartheta(\mathrm x)).
\end{aligned}
\]

\uline{Exemple} 1.- Une variété de groupe est un noyau de groupe.

\uline{Exemple} 2.- Soient $\mathrm X$ une variété, $\mathrm e\in\mathrm X$ et $\mathrm m$ une application d'un ouvert $\Omega$ de $\mathrm X\times\mathrm X$ dans $\mathrm X$ satisfaisant aux conditions (ii) et (iii) de la définition 5. Le lemme 1 du n\textsuperscript{o} 3 montre qu'il existe un voisinage ouvert $\mathrm G$ de $\mathrm e$ dans $\mathrm X$ et un automorphisme $\vartheta$ de $\mathrm G$ sur lui-même satisfaisant à (i). Il suffit alors de prendre pour $\Omega(\mathrm G)$ l'ensemble des $(\mathrm x,\mathrm y)\in(\mathrm G\times\mathrm G)\cap\Omega$ tels que $\mathrm m(\mathrm x,\mathrm y)\in\mathrm G$ et pour $\mathrm m_{\mathrm G}$ la restriction de $\mathrm m$ à $\Omega(\mathrm G)$ pour obtenir sur $\mathrm G$ une structure de noyau de groupe.
% archive sheet number: 215

\scanpage{217}
Lien avec la variété des modules pour des \unsure{multip.|multips|motifs} polarisées rigidifiées : Si \deleted{\unsure{$A$|$S$}} $\unsure{\mathcal M|\mathscr M}$ est une telle variété modulaire, alors\unsure{,} un morphisme $X\unsure{\mathrel{\text{---}}|\longrightarrow}\mathcal M$ est déterminé quand on connaît sa valeur \uline{en un point}, plus l'application $\pi_1(x)\mathrel{\text{---}}\pi_1(\mathcal M)$ ......

C'est une propriété intrinsèque importante de ces $\mathcal M$, et une propriété de cette nature mérite d'être étudiée pour \unsure{ses|les} propres mérites .... Peut-être \unsure{peuvent|pourront} elles \unsure{ttes|toutes} être \deleted{\unsure{considérées}} \inserted{interprétées} comme des \deleted{\unsure{schémas|espaces}} \inserted{variétés} modulaires pour certains types de motifs.

\begin{center}\rule{3cm}{0.5pt}\end{center}

Introduire la notion de schéma \inserted{motiviquement} \unsure{ponctué|pointé} \inserted{\tikz[baseline=(a.base)]\node[draw,ellipse,inner sep=1pt](a){(à val. dans un anneau $k$)};}, et morphismes de tels, qui définissent des morphismes correspondants pour les
\[
\begin{aligned}
&G(X,\unsure{\xi|\zeta})\mathrel{\text{---}}G(Y,\unsure{\eta|y}),\qquad\text{donc pour}\\
&G(X,\xi)(k)=\pi_1(X,\xi)\mathrel{\text{---}}G(Y,\eta)(k)=\pi_1(Y,\eta).
\end{aligned}
\]
% archive sheet number: 216

